%=========================================================================
% Memory in Memory: Learning Procedural Memory Policies
% from Downstream Failures
%
% Status: draft of Abstract + Introduction + Related Work + Method.
% The Experiments section (Section 5) is intentionally not written yet.
%=========================================================================
\documentclass[11pt]{article}

%-------------------------------------------------------------------------
% Packages
%-------------------------------------------------------------------------
\usepackage[margin=1in]{geometry}
\usepackage{amsmath,amssymb,amsthm}
\usepackage{mathtools}
\usepackage{graphicx}
\usepackage{booktabs}
\usepackage{xcolor}
\usepackage{enumitem}
\usepackage{algorithm}
\usepackage{algpseudocode}
\usepackage{microtype}
\usepackage[hidelinks]{hyperref}

%-------------------------------------------------------------------------
% Theorems and custom environments
%-------------------------------------------------------------------------
\newtheorem{definition}{Definition}[section]
\newtheorem{remark}{Remark}[section]
\newtheorem{assumption}{Assumption}[section]

%-------------------------------------------------------------------------
% Math macros
%-------------------------------------------------------------------------
\newcommand{\MM}{\mathcal{M}}
\newcommand{\DD}{\mathcal{D}}
\newcommand{\KK}{\mathcal{K}}
\newcommand{\SSs}{\mathcal{S}}
\newcommand{\CC}{\mathcal{C}}
\newcommand{\RR}{\mathcal{R}}
\newcommand{\EE}{\mathbb{E}}
\newcommand{\II}{\mathbb{I}}
\DeclareMathOperator*{\argmax}{arg\,max}
\DeclareMathOperator{\TopK}{TopK}
\DeclareMathOperator{\Cover}{Cover}
\DeclareMathOperator{\Support}{Support}
\DeclareMathOperator{\Apply}{Apply}
\DeclareMathOperator{\Valid}{Valid}
\DeclareMathOperator{\Cluster}{Cluster}
\DeclareMathOperator{\Rank}{rank}
\DeclareMathOperator{\simsem}{sim_{sem}}
\DeclareMathOperator{\simlex}{sim_{lex}}
\DeclareMathOperator{\Cand}{G_{cand}}
\DeclareMathOperator{\Crud}{G_{crud}}

\title{\textbf{Memory in Memory: Learning Procedural Memory Policies from Downstream Failures}}

\author{%
  Anonymous Author(s) \\
  \small Contact information omitted for anonymization
}

\date{August 2026}

\begin{document}

\maketitle

%=========================================================================
% ABSTRACT
%=========================================================================
\begin{abstract}
\noindent
Large language model (LLM) agents increasingly rely on external memory to preserve
user facts, state, and history across sessions. Existing memory systems primarily
improve the \emph{representation}, \emph{compression}, and \emph{retrieval} of
content, but rarely study how a memory system can continuously learn
``how it should remember'' from its own downstream errors. When a question is
answered incorrectly, the same failure signal can originate from completely
different stages of the pipeline: the target information may never have been
written into memory, it may exist but fail to be retrieved, or it may already be
in context yet be misused by the answer model. Methods that reflect over failures
or rewrite global prompts without separating these root causes risk misattribution,
duplicated rules, and global regressions.

We propose \textbf{Memory in Memory} (MiM), an error-driven \emph{procedural
meta-memory} framework for LLM-mediated memory systems. MiM decomposes a long-term
memory system into a memory construction policy and a memory access policy, and
adds an offline error-learning loop on top of normal operation. Given a downstream
failure, MiM runs three information-isolated checks---answer evidence sufficiency,
access coverage of currently useful memory, and construction fidelity along
provenance and version chains---so that access omissions and construction losses
are attributed as non-exclusive, causally scoped failures rather than a single
miscellaneous label. Confirmed failures are distilled into natural-language
\emph{candidate skills}, which are semantically clustered, related to the existing
skill bank through an explicit candidate--bank affinity matrix, and integrated by
a constrained batch CRUD process with write-set conflict resolution before being
released as a new versioned skill bank. At runtime, only frozen official skills
are retrieved, while a skill-exposure trace records both the skills actually
injected and the boundary candidates near the retrieval margin, enabling later
learning rounds to distinguish policy absence, recall failure, and execution
ineffectiveness. MiM thereby lets a long-term memory system accumulate not only
object-level facts about users and the world, but also procedural experience
about how its own construction and access policies should behave.
\end{abstract}

% TODO(experiments): once Section 5 is written, add a sentence reporting the
% headline empirical results (e.g., C+P rate and token F1 under the frozen
% protocol) to close the abstract.

%-------------------------------------------------------------------------
% KEYWORDS (uncomment if the venue requires them)
%-------------------------------------------------------------------------
% \noindent\textbf{Keywords:} agent memory, procedural memory, failure attribution,
% skill learning, LLM agents

%=========================================================================
% 1. INTRODUCTION
%=========================================================================
\section{Introduction}
\label{sec:intro}

External memory has become a foundational component of long-running LLM
agents~\cite{packer2023memgpt,du2026survey}. An agent equipped with persistent
memory can retain user preferences, personal relationships, historical events,
and state changes across sessions, and selectively recover the relevant
information in future tasks~\cite{park2023generative,wu2024longmemeval,
maharana2024locomo}. Toward this goal, a rich design space has emerged: vector
memory, hierarchical summaries, structured facts, temporal knowledge graphs,
multi-type memory, and agent-managed memory systems~\cite{packer2023memgpt,
zhang2025amem,chattopadhyay2025mem0,chen2025memoryos,vlasic2025mirix,zhu2025graphiti}.
Despite ever-stronger underlying representations, long-term memory systems still
exhibit recurring failure patterns: important information is not retained,
updates destroy historical state, semantically similar but distinct events are
incorrectly merged, a fact that exists in memory is never retrieved, or a
multi-hop question is answered prematurely with only partial evidence.

These errors do not originate solely from storage capacity or retriever
performance. For an \emph{LLM-mediated} memory system, the decisive decisions
occur at the policy level: the model decides which information is worth writing,
how new information should be integrated with existing memory, what should be
retrieved for the current question, whether to keep searching, and when the
accumulated evidence is sufficient to answer. Even if the underlying store can
persist and return data faithfully, inappropriate construction and access
policies produce systematic failures.

Existing systems typically accumulate \emph{object-level} memory, that is, facts
about users, environments, and history. In contrast, they rarely accumulate
\emph{procedural meta-memory}, i.e., experience about the memory system's own
operation. For example, a system may record ``the user has moved to Shanghai''
but never learn from a previous mistake that ``when updating the current address,
the historical validity interval of the old address should be preserved''; it
may store several related events but never learn that ``for list questions, one
should not stop retrieving after finding only a single component.''

A straightforward remedy is to save failure cases and retrieve past errors when a
similar problem arises, or to append every reflection to the global prompt.
Both approaches lack a critical \emph{root-cause constraint}. An incorrect answer
is only a downstream symptom: the target information may have failed to enter
memory, failed to be retrieved, or been visible yet misused. If the system
generalizes an access omission into a construction rule, or writes an answer
reasoning mistake as a retrieval policy, the new experience will not only fail to
transfer but may also harm unrelated tasks.

Learning memory policies from errors therefore first requires answering a causal
question:

\begin{quote}
\emph{In the pipeline spanned by raw interaction, memory state, access trajectory,
and answer context, where is the first point at which the target information is
no longer correctly preserved or used?}
\end{quote}

We propose \textbf{Memory in Memory} (MiM), which turns downstream errors into
attributable, maintainable, and reusable procedural memory policies. MiM does not
replace the underlying memory system; it adds an external \emph{control plane}.
During normal operation, the base system constructs and accesses memory exactly
as before. In an offline learning phase, MiM observes the raw interaction, memory
versions, natural search trajectories, and answer outputs, and diagnoses three
families of failure: evidence misuse at answer time, access omission over the
current memory state, and construction loss along the version history.

MiM's diagnosis is not a single mutually exclusive label produced by one model
call. Answer sufficiency, access omission, and construction loss are modeled as
three checks over \emph{disjoint information views}. The access check compares
the currently useful memory and the natural search results on the fixed current
snapshot, without reading the raw conversation or historical versions; the
construction check first verifies whether the current memory is complete and
only then walks the source and version relations to locate the earliest
corruption. A single question can carry both access and construction problems,
because the current memory may have lost one fact while simultaneously failing
to recall another fact that is still present.

For confirmed access or construction failures, MiM does not inject the full
failure case into the running context. Instead it distills a short, retrievable,
natural-language \emph{candidate skill} that describes when a class of situations
occurs and what operational principle the memory agent should follow. To prevent
per-case updates from inflating the policy library, MiM first induces candidates
independently, then clusters semantically similar candidates, retrieves the
related official skills, and executes a constrained batch CRUD consolidation.
Only after global integration and consistency validation does a candidate become
an official, runtime-retrievable policy.

MiM also records, for every decision, which official skills were actually
retrieved and which candidates sat near the retrieval boundary. This
\emph{skill-exposure trace} lets the next learning round distinguish three
different failure types: no relevant experience exists in the policy bank, a
relevant experience exists but did not enter the context, and a relevant
experience was used but its content was ineffective. Skill learning thus ceases
to be an unobservable prompt rewrite and becomes a closed loop whose
representation, recall, and execution stages can be analyzed separately.

Our core contributions are as follows:

\begin{enumerate}[leftmargin=*,itemsep=2pt,topsep=4pt]
  \item We introduce the problem setting of \emph{error-driven procedural
  meta-memory}: long-term memory systems should accumulate reusable experience
  about their own construction and access policies, not only object-level facts.
  \item We propose an \emph{information-isolated failure attribution} method that
  decomposes answer insufficiency, access omission, and construction loss into
  non-exclusive checks, and uses source and version relations to locate the
  first state transition at which information fidelity is broken.
  \item We propose \emph{direction-separated dual skill banks} for access and
  construction, together with a two-stage promotion mechanism from
  instance-level candidates to official frozen policies.
  \item We propose a \emph{constrained skill-bank evolution} procedure based on
  candidate--bank relations, batch CRUD with a deterministic validator, and
  write-set conflict resolution, in order to control duplication, order
  dependence, and version oscillation.
  \item We build a fully observable long-term-memory research base and define a
  conversation-level frozen evaluation protocol that measures final task quality,
  error-type changes, and natural skill invocation separately.
\end{enumerate}

The remainder of this paper is organized as follows. Section~\ref{sec:related}
discusses related work on agent memory, self-evolving memory policies, and
error-driven agent improvement. Section~\ref{sec:problem} formalizes
LLM-mediated memory systems, object-level versus procedural memory, and the
failure signal. Section~\ref{sec:method} presents the MiM framework in detail,
including traceable memory states, closed-loop access, isolated failure
attribution, candidate induction, semantic consolidation, and constrained bank
evolution. Section~\ref{sec:conclusion} concludes.

%=========================================================================
% 2. RELATED WORK
%=========================================================================
\section{Related Work}
\label{sec:related}

\subsection{Memory Architectures for LLM Agents}
\label{sec:rw-arch}

Persistent memory for LLM agents has been studied along several axes. \textbf{Tiered
memory} organizes information by access frequency and lifetime: MemGPT
(Letta)~\cite{packer2023memgpt} borrows the operating-system metaphor of
working memory and long-term storage and exposes memory manipulation as function
calls, making access decisions part of the agent's tool-use loop. \textbf{Structured
and associative memory} represents facts as knowledge graphs or associative
networks, as in Mem0~\cite{chattopadhyay2025mem0}, A-MEM~\cite{zhang2025amem},
and Zep/Graphiti~\cite{zhu2025graphiti}, which build temporal knowledge graphs
that support time-aware updates. \textbf{Multi-type memory} distinguishes
episodic, semantic, procedural, and other categories: MemoryOS~\cite{chen2025memoryos}
layers short-term, mid-term, and long-term stores; MIRIX~\cite{vlasic2025mirix}
routes inputs through six specialized memory managers supervised by a meta
manager and pairs them with an access agent that selects memory components and
retrieval methods per question. \textbf{Efficient memory} targets token economy:
LightMem~\cite{luo2025lightmem} and SimpleMem~\cite{liu2026simplemem} compress
interaction histories into retrieval-friendly units, and Hindsight~\cite{van2025hindsight}
maintains four typed networks with retain/recall/reflect operations. Recent
surveys~\cite{du2026survey} formalize agent memory as a write--manage--read
loop, a perspective that positions policy control---what to write, how to manage,
and how to read---as the core design decision.

These systems improve \emph{what} is stored and \emph{how} it is retrieved, but
the policies that govern construction and access are largely fixed at deployment.
MiM complements this line by treating construction and access \emph{policies}
themselves as a learnable resource: rather than proposing a new storage
representation, it attaches an error-learning loop that accumulates procedural
experience about any underlying memory architecture.

\subsection{Skill-Guided and Self-Evolving Memory Policies}
\label{sec:rw-self}

A recent line of work makes memory policies learnable. \textbf{MemSkill}~\cite{zhang2026memskill}
reframes memory construction operations as a small set of reusable, evolvable
skills: a controller selects relevant skills, an executor produces skill-guided
memories, and a designer periodically reviews hard cases to refine or create
skills. MemSkill is the closest neighbor on the construction side, but its skill
selection is trained with reinforcement learning and it does not learn a
multi-tool, stop-capable access policy. \textbf{EvolveMem}~\cite{liu2026evolvemem}
exposes the full retrieval configuration (scoring functions, fusion weights,
budgets) as an action space and evolves it with an LLM-powered diagnosis module
under revert-on-regression guards. It targets continuous configuration
optimization rather than discrete, retrievable, natural-language policies.
\textbf{MemMA}~\cite{memma2026} places a meta-thinker and query reasoner outside
a unified storage interface: it diagnoses evidence gaps per question, generates
orthogonal queries, and directly repairs factual entries on the current instance.
MemMA does not accumulate cross-instance procedural skills; its controller
improves the current question rather than the agent's future construction and
access behavior. \textbf{PlugMem}~\cite{plugmem2026} is a task-agnostic plugin
controller for memory access, selecting memory types and planning retrieval, but
it is designed as a replaceable module rather than a bank of experience learned
from failures. On the access side, \textbf{MRAgent}~\cite{ji2026mragent} and
\textbf{MemMachine}~\cite{memmachine2026} build strong multi-tool retrieval
agents over graph or episode stores, yet their strategies are frozen after
design.

MiM differs from this line in two ways. First, it is \emph{error-driven and
attribution-first}: nothing is learned unless a downstream failure has been
isolated to the access or construction stage, so repair targets are causally
scoped. Second, its skills are \emph{procedural meta-memory}: versioned,
retrievable natural-language policies that are integrated through constrained
batch CRUD, rather than continuous configuration deltas or per-instance fact
repairs.

\subsection{Error-Driven Agent Improvement and Reflection}
\label{sec:rw-reflect}

Reflection-based methods let an agent improve from its own mistakes.
Reflexion~\cite{shinn2023reflexion} converts failed trajectories into verbal
feedback stored in an episodic buffer; Self-Refine~\cite{madaan2023selfrefine}
iteratively revises outputs; Generative Agents~\cite{park2023generative}
maintain episodic buffers indexed by recency and importance; MemoryBank~\cite{zhong2024memorybank}
applies forgetting curves and periodic consolidation to prune stale entries.
Prompt-level optimization, such as OPRO~\cite{yang2024opro} and the prompt
optimizer shipped with LangMem~\cite{langmem2025}, rewrites the agent's policy
prompt globally from trajectory feedback.

These methods share a common weakness that MiM targets directly: reflection is
applied \emph{without separating where the pipeline first failed}. A global
prompt rewrite triggered by an access failure may modify construction behavior;
a reflection over an answer reasoning error may produce a retrieval rule. MiM
instead attributes each failure to answer, access, or construction stages first,
and only then induces a directional skill that is retrieved conditionally at
runtime, so new experience is applied where its root cause lives and remains
reusable without contaminating unrelated tasks.

\subsection{Skill Libraries for Agents}
\label{sec:rw-skill}

Beyond memory, agents increasingly reuse \emph{procedural skills}. Voyager~\cite{wang2023voyager}
maintains a growing library of executable code skills acquired through
environment feedback, and recent agent frameworks treat skills as versioned,
retrievable units of behavior. MiM shares the spirit of \emph{learning reusable
procedures rather than facts}, but with a different substrate: its skills are
concise natural-language operation principles for memory construction and
access, induced from offline failure attribution, and maintained through a
deterministically validated, conflict-free batch evolution process. This makes
skill quality inspectable and the bank versionable, both of which are
prerequisites for a memory system whose learning must not compromise the
reliability of the object-level store it guides.

%=========================================================================
% 3. PROBLEM FORMULATION
%=========================================================================
\section{Problem Formulation}
\label{sec:problem}

\subsection{LLM-Mediated Memory Systems}
\label{sec:prob-sys}

We consider a long-term memory system whose construction and access decisions
are at least partially made by LLMs:

\begin{equation}
\label{eq:system}
\MM \;=\; \left(\DD,\;\pi_{\mathrm{con}},\;\pi_{\mathrm{acc}},\;f_{\mathrm{ans}}\right),
\end{equation}

\noindent where $\DD$ is the persistent memory state, $\pi_{\mathrm{con}}$ is the
memory \emph{construction policy}, $\pi_{\mathrm{acc}}$ is the memory
\emph{access policy}, and $f_{\mathrm{ans}}$ is the task model that produces an
answer from the question and the visible evidence.

Given the $t$-th interaction segment $x_t$ and the previous memory state
$D_{t-1}$, the construction policy produces a new state and a construction
trajectory:

\begin{equation}
\label{eq:construction}
D_t,\;\tau_t^{\mathrm{con}} \;=\; \pi_{\mathrm{con}}\left(x_t,\,D_{t-1};\;S_t^{\mathrm{con}}\right),
\end{equation}

\noindent where $S_t^{\mathrm{con}}$ is the set of construction skills available
at this round and $\tau_t^{\mathrm{con}}$ is the construction trajectory (the
sequence of candidate proposals, decisions, and state transitions).

Given a question $q$, the access policy interacts with memory through a sequence
of actions $a_1,\dots,a_J$:

\begin{equation}
\label{eq:access}
a_j \;\sim\; \pi_{\mathrm{acc}}\left(a_j \mid q,\,D_t,\,S_q^{\mathrm{acc}},\,h_{<j}\right),
\end{equation}

\noindent where $S_q^{\mathrm{acc}}$ is the set of access skills retrieved for
$q$, and $h_{<j}$ contains the previous queries, tool observations, and
intermediate judgments. When the policy executes a stop action, the answer model
produces a prediction from the accumulated visible evidence
$M_{\mathrm{visible}}$:

\begin{equation}
\label{eq:answer}
\hat{y} \;=\; f_{\mathrm{ans}}\left(q,\;M_{\mathrm{visible}}\right).
\end{equation}

In our current implementation, access and answering are performed by the same
closed-loop agent, so the stop decision and the final answer share the full
natural search context. We revisit this choice in
Section~\ref{sec:method-rationale}.

\subsection{Object-Level and Procedural Memory}
\label{sec:prob-object}

Object-level memory $D_t$ stores facts about users and the world. MiM maintains
an additional, separate resource: two \emph{procedural policy banks}

\begin{equation}
\label{eq:banks}
\KK \;=\; \left(\KK_{\mathrm{con}},\;\KK_{\mathrm{acc}}\right),
\end{equation}

\noindent where $\KK_{\mathrm{con}}$ guides information selection, extraction,
updating, merging, and historical preservation during construction, and
$\KK_{\mathrm{acc}}$ guides query planning, retrieval routes, filtering,
evidence aggregation, and stopping during access.

Each skill uses a minimal natural-language representation

\begin{equation}
\label{eq:skill}
s \;=\; (n,\,d,\,c),
\end{equation}

\noindent where $n$ is the skill name, $d$ is the retrieval description
(``when does this skill apply''), and $c$ is the execution content (``what
principle should be followed''). Skills do not retain user facts from training
cases; they describe processing principles that transfer across entities and
sessions.

\subsection{Failure Signal}
\label{sec:prob-failure}

Given a reference answer $y^{*}$ and the model prediction $\hat{y}$, a semantic
judge yields

\begin{equation}
\label{eq:judge}
J(\hat{y},\,y^{*}) \;\in\; \{\textsc{C},\,\textsc{P},\,\textsc{I}\},
\end{equation}

\noindent where \textsc{C}, \textsc{P}, and \textsc{I} denote correct, partial,
and incorrect, respectively. Only \textsc{P} and \textsc{I} enter failure
attribution. The judge observes only the question, the reference answer, the
prediction, and the conversation time anchor; it does not read the memory state
or the search trajectory, so it confirms a downstream failure without explaining
it.

Let $\Cover(M,\,y^{*})$ denote the degree to which an evidence set $M$ covers the
necessary facts of $y^{*}$, and $\Support(m,\,y^{*})$ indicate whether a memory
entry $m$ supports at least one necessary fact. MiM defines three diagnosis
outcomes.

\begin{definition}[Answer failure]
\label{def:ans}
An \emph{answer failure} occurs when the prediction is not correct even though
the evidence visible to the answer model was already sufficient:
\begin{equation}
\label{eq:fans}
F_{\mathrm{ans}} \;=\; \II\left[\,J(\hat{y},y^{*})\neq \textsc{C}
\;\land\; \Cover(M_{\mathrm{visible}}, y^{*}) = 1\,\right].
\end{equation}
\end{definition}

\begin{definition}[Access failure]
\label{def:acc}
Let the currently useful memory be
\begin{equation}
\label{eq:museful}
M_{\mathrm{useful}} \;=\; \left\{\, m \in D_t \;:\;
\Support(m,\,y^{*}) = 1 \,\right\},
\end{equation}
\noindent and let the memories returned by the natural search chain be
\begin{equation}
\label{eq:mretrieved}
M_{\mathrm{retrieved}} \;=\; \bigcup_{j=1}^{J} M_j .
\end{equation}
\noindent An \emph{access failure} occurs when useful memory exists but was not
retrieved:
\begin{equation}
\label{eq:facc}
F_{\mathrm{acc}} \;=\; \II\left[\;
M_{\mathrm{useful}} \setminus M_{\mathrm{retrieved}} \;\neq\; \varnothing
\,\right].
\end{equation}
\end{definition}

\begin{definition}[Construction failure]
\label{def:con}
A \emph{construction failure} occurs when the raw interaction supports the
necessary facts but the current memory state no longer covers them:
\begin{equation}
\label{eq:fcon}
F_{\mathrm{con}} \;=\; \II\left[\;
\Support(X_{\mathrm{raw}},\,y^{*}) = 1
\;\land\; \Cover(D_t,\,y^{*}) < 1 \,\right],
\end{equation}
\noindent where $X_{\mathrm{raw}}$ denotes the raw source evidence.
\end{definition}

These outcomes are not mutually exclusive. In particular,
$F_{\mathrm{acc}}$ and $F_{\mathrm{con}}$ can hold simultaneously: the current
memory may have lost one fact (construction) while failing to recall another
fact that is still present (access). MiM therefore never reduces a failure to a
single label; the three checks are computed independently over disjoint
information views (Section~\ref{sec:method-attribution}).

%=========================================================================
% 4. METHOD
%=========================================================================
\section{Method: Memory in Memory}
\label{sec:method}

\subsection{Overview}
\label{sec:method-overview}

MiM adds an error-learning control plane outside the base memory system. One
complete iteration consists of three phases:

\begin{enumerate}[leftmargin=*,itemsep=2pt,topsep=4pt]
  \item the base system constructs memory and answers training questions under a
  \emph{frozen} skill bank;
  \item MiM runs information-isolated root-cause diagnosis over semantic
  failures and induces candidate skills;
  \item candidates are globally consolidated into a new skill-bank version,
  which is selected on validation data and then frozen for subsequent runs.
\end{enumerate}

This design separates three concerns that are usually conflated. \emph{Diagnosis}
answers ``what does this failure require learning''; \emph{consolidation}
answers ``how does this repair relate to existing policies''; \emph{frozen
evaluation} answers ``does the updated bank work on unseen sessions.''
Algorithm~\ref{alg:train} summarizes the training-time loop.

\begin{algorithm}[t]
\caption{Failure-driven skill learning (training phase)}
\label{alg:train}
\begin{algorithmic}[1]
\Require conversations, questions, train split, frozen bank $\KK^{(v)}$
\State \textbf{Phase 1 -- run.}
\For{conversation $c$ in train split}
  \State $D_t \leftarrow$ \textsc{Ingest}($c$;\; $\KK_{\mathrm{con}}^{(v)}$) \Comment{batch CRUD, single transaction}
  \For{question $q$ in $c$}
    \State $(\hat{y}, M_{\mathrm{visible}}, \tau^{\mathrm{acc}}) \leftarrow$ \textsc{Ask}($q$;\; $\KK_{\mathrm{acc}}^{(v)}$)
    \State record skill-exposure trace $T_q = (\mathcal{S}^{\mathrm{selected}}_q, \mathcal{S}^{\mathrm{nearby}}_q, v)$
  \EndFor
\EndFor
\State \textbf{Phase 2 -- attribute.}
\For{failed question with $J(\hat{y},y^{*})\neq\textsc{C}$}
  \State $F_{\mathrm{ans}}$: answer check on $M_{\mathrm{visible}}$ only
  \State $F_{\mathrm{acc}}$: useful vs.\ retrieved on current snapshot $D_t$
  \State $F_{\mathrm{con}}$: screen $D_t$; if gap found, trace first error $t^{*}$
  \For{side $z \in \{\mathrm{acc},\mathrm{con}\}$ with confirmed failure}
    \State $c_i \leftarrow \Cand(d_i, T_i)$ \Comment{or $c_i = \varnothing$ (no change)}
  \EndFor
\EndFor
\State \textbf{Phase 3 -- consolidate and publish.}
\State $(\KK_{\mathrm{acc}}^{(v+1)}, \KK_{\mathrm{con}}^{(v+1)}) \leftarrow$ \Call{Consolidate}{candidates, $\KK^{(v)}$} \Comment{Alg.~\ref{alg:consolidate}}
\State select bank version on validation split; freeze for evaluation
\end{algorithmic}
\end{algorithm}

\subsection{Traceable Memory States}
\label{sec:method-states}

Attributing construction failures requires that the memory system expose source
and version relations. We represent each memory state as

\begin{equation}
\label{eq:state}
D_t \;=\; \left(V_t,\; E_{\mathrm{src}},\; E_{\mathrm{par}},\; E_{\mathrm{chg}}\right),
\end{equation}

\noindent where $V_t$ is the set of memory versions visible at system time $t$,
$E_{\mathrm{src}}$ connects versions to their originating raw messages,
$E_{\mathrm{par}}$ denotes parent--child relations introduced by updates or
merges, and $E_{\mathrm{chg}}$ records the state transitions that produced new
versions. Each transition is generated by an operation

\begin{equation}
\label{eq:ops}
o_t \;\in\; \{\textsc{Add},\,\textsc{Update},\,\textsc{Merge},\,\textsc{Delete},\,\textsc{Skip}\},
\end{equation}

\noindent and source relations are inherited across updates, so the system can
trace any raw message to the candidate it first formed, its initial version, and
every subsequent change.

This representation lifts construction diagnosis from ``is the final memory
wrong'' to ``at which state transition was the information first corrupted''.
Let $\mathcal{I}(D_t, X)$ be the invariant that memory state $D_t$ still
faithfully preserves the target information from source evidence $X$. The
earliest error point is

\begin{equation}
\label{eq:tstar}
t^{*} \;=\; \min\left\{\, t \;:\;
\mathcal{I}(D_{t-1}, X) = 1 \;\land\; \mathcal{I}(D_t, X) = 0 \,\right\}.
\end{equation}

\noindent If the initial extraction was already wrong, $t^{*}$ corresponds to
candidate formation; if the initial state was correct, $t^{*}$ corresponds to
the first update, merge, or deletion that broke the invariant.

\subsection{Closed-Loop Memory Access}
\label{sec:method-access}

MiM does not model access as a single fixed retrieval. The access policy
chooses, after every observation, between searching, inspecting, and answering.
At round $j$, the policy receives the previous history $h_{<j}$, emits action
$a_j$, receives tool observation $o_j$, and updates

\begin{equation}
\label{eq:hist}
h_j \;=\; h_{j-1} \cup \{a_j,\,o_j\}.
\end{equation}

Search actions can choose among semantic, lexical, exact-match, and structured
temporal routes. For hybrid retrieval, the candidate score is

\begin{equation}
\label{eq:hybrid}
s(m \mid q) \;=\; \sum_{r \in \RR} \frac{w_r}{k + \Rank_r(m)},
\end{equation}

\noindent where $\RR$ is the set of retrieval routes with weights $w_r$ summing
to one (reciprocal-rank fusion), and is then corrected by entity matching, target
time validity, and current-state activity factors.

The agent judges, from the accumulated evidence, whether the necessary facts are
complete. If information is insufficient, it can decompose the question, change
the query formulation, switch routes, or deepen the search; if sufficient, it
stops immediately. This closed loop gives access skills a real control target:
a skill can influence not only the final query text but also the retrieval
route, the continuation condition, and the evidence-completeness judgment.

\subsection{Isolated Failure Attribution}
\label{sec:method-attribution}

MiM builds three independent information views for the same failure.

\paragraph{Answer view.}
The answer check sees only the evidence actually visible to the running model,
and determines whether the error occurred at the evidence-utilization stage
(Definition~\ref{def:ans}). It produces no repair package and does not enter
skill learning; it is recorded so that failures are never misattributed to
memory when the memory evidence was already sufficient.

\paragraph{Access view.}
The access check is fixed to the current memory snapshot and contains only the
currently relevant memory and the natural access trajectory. It does not read
the raw conversation or historical versions, so that upstream construction
errors are not introduced into access attribution. A semantic model labels the
useful entries, and a deterministic program computes the difference

\begin{equation}
\label{eq:missing}
M_{\mathrm{missing}} \;=\; M_{\mathrm{useful}} \setminus M_{\mathrm{retrieved}}.
\end{equation}

\paragraph{Construction view.}
The construction check uses progressive disclosure. The first stage examines
only whether the current relevant memory covers the necessary facts
(\textsc{Full}/\textsc{Partial}/\textsc{Missing}/\textsc{Incorrect}); only when
a problem is found does the second stage open the raw source messages and the
version trajectory to locate $t^{*}$ via Eq.~\eqref{eq:tstar}.

This permission design is not about shrinking model inputs; it establishes an
explanatory causal boundary. An access failure means ``available information was
not retrieved,'' while a construction failure means ``the current information
state itself is wrong.'' Because the three workflows are physically and
contextually isolated, they can run in parallel and produce independent repair
targets.

\subsection{Runtime Skill Exposure}
\label{sec:method-exposure}

At runtime, the agent retrieves skills from the direction-corresponding bank:

\begin{equation}
\label{eq:skillretr}
\mathcal{S}_q \;=\; \TopK_{s \in \KK_z}\left[\;
\lambda \cdot \simsem(q, s) \;+\; (1-\lambda) \cdot \simlex(q, s)
\,\right], \qquad z \in \{\mathrm{acc},\,\mathrm{con}\},
\end{equation}

\noindent where semantic similarity is computed over the skill name and
description (the trigger condition), and lexical similarity is a BM25 score.
A fixed candidate pool is drawn first, then an LLM-based applicability reranker
filters out skills that do not actually apply to the current decision, and a
score threshold rejects low-confidence matches. In our configuration,
$\lambda = 0.7$, at most two skills are injected per decision, and the five
next-ranked boundary candidates are retained.

Besides the injected skills, the system records the candidates near the
retrieval boundary that did not enter the context:

\begin{equation}
\label{eq:trace}
T_q \;=\; \left(\mathcal{S}_q^{\mathrm{selected}},\;
\mathcal{S}_q^{\mathrm{nearby}},\; v_{\KK}\right),
\end{equation}

\noindent where $v_{\KK}$ is the official bank version. $T_q$ is the policy
exposure trace consumed by later learning rounds. If a relevant skill already
entered $\mathcal{S}_q^{\mathrm{selected}}$ and the failure persists, the
problem most likely lies in the skill content or its execution; if it appears
only in $\mathcal{S}_q^{\mathrm{nearby}}$, the problem is likely in the retrieval
representation; if it appears in neither, the system may need to learn a new
policy. These three cases---policy absence, recall failure, and execution
ineffectiveness---are thereby separable, which is precisely what a
prompt-level reflection method cannot observe.

\subsection{Failure-Conditioned Skill Induction}
\label{sec:method-induction}

For each diagnosis package $d_i$ and its policy exposure trace $T_i$, a
candidate agent produces

\begin{equation}
\label{eq:candidate}
c_i \;=\; \Cand(d_i,\,T_i),
\end{equation}

\noindent where $c_i$ contains a name, a retrieval description, execution
content, and a short statement of the general problem it solves. The generation
process is required to remove concrete entities, dates, answers, and internal
IDs so that the candidate remains reusable across cases.

The candidate agent may also output

\begin{equation}
\label{eq:nochange}
c_i \;=\; \varnothing,
\end{equation}

\noindent indicating that existing policies already cover the problem, or that
the problem is not suitable for repair through a memory skill. The empty
operation is a deliberate mechanism for controlling bank growth, not a
generation failure.

Candidates cannot be retrieved by the runtime before publication. MiM therefore
establishes a two-stage policy lifecycle: instance-level failures first produce
non-executable repair hypotheses, and a global consolidation process decides
their creation, merging, update, or rejection.

\subsection{Semantic Candidate Consolidation}
\label{sec:method-consolidation}

Updating the policy bank error-by-error introduces strong order dependence:
a skill released for one error changes the bank seen by the next error, similar
failures may create duplicate rules, and the same skill may be rewritten
repeatedly. MiM therefore collects candidates first and integrates them in
batches.

The semantic representation of a candidate is

\begin{equation}
\label{eq:cembed}
e(c_i) \;=\; \alpha\, e_{\mathrm{desc}}(c_i) + \beta\, e_{\mathrm{content}}(c_i)
+ \gamma\, e_{\mathrm{solves}}(c_i),
\end{equation}

\noindent with $\alpha + \beta + \gamma = 1$ (in our configuration,
$0.45 / 0.35 / 0.20$). Semantically similar candidates are grouped by
deterministic spherical $K$-means into candidate clusters
$\{\CC_g\}_{g=1}^{G}$, subject to a target cluster size and a maximum batch size.

For a candidate $c_i$ and an existing skill $s_j$, MiM computes the relation
explicitly:

\begin{equation}
\label{eq:relation}
R_{ij} \;=\; \lambda_d\, \simsem\!\left(c_i^{d},\, s_j^{d}\right)
\;+\; \lambda_c\, \simsem\!\left(c_i^{c},\, s_j^{c}\right)
\;+\; \lambda_l\, \simlex(c_i, s_j),
\end{equation}

\noindent with $\lambda_d + \lambda_c + \lambda_l = 1$ (in our configuration,
$0.50 / 0.30 / 0.20$). Compared with concatenating an entire group of candidates
into a single query, the matrix $R$ preserves each repair hypothesis's own
related skills while still allowing the discovery of a common policy that covers
multiple candidates. Before planning, candidates whose best relation score
exceeds a threshold (0.60 in our configuration) are linked to the corresponding
official skills.

\subsection{Constrained Skill-Bank Evolution}
\label{sec:method-evolution}

For a candidate cluster $\CC_g$, the related existing skills, and the relation
matrix $R_g$, a CRUD agent produces an update plan

\begin{equation}
\label{eq:crud}
P_g \;=\; \Crud\left(\CC_g,\; \KK_z^{(v)},\; R_g\right).
\end{equation}

\noindent The plan can add a skill, rename a skill, update a description, add
or update content, move content, or deactivate a skill. Every candidate must
receive a unique resolution, e.g., created, merged into an existing skill,
covered by an existing skill, or rejected. The model only proposes plans; a
deterministic validator checks target existence, direction, expected bank
version, content positions, and candidate coverage:

\begin{equation}
\label{eq:valid}
\Valid(P_g,\; \KK_z^{(v)}) = 1 .
\end{equation}

\noindent Only valid plans may modify the bank.

Different candidate clusters can be planned in parallel against the same frozen
version. Let $W_g$ be the write set of plan $P_g$ (the set of skills it
modifies). If

\begin{equation}
\label{eq:conflict}
W_i \cap W_j \neq \varnothing,
\end{equation}

\noindent the two plans conflict. MiM builds the conflict graph and re-generates
a joint plan for each connected component, re-retrieving the combined candidate
set against the official bank. After all conflicts are resolved, the updates are
applied in a single transaction to produce a new bank version:

\begin{equation}
\label{eq:release}
\KK_z^{(v+1)} \;=\; \Apply\left(\KK_z^{(v)},\; \bigcup_g P_g\right).
\end{equation}

\noindent This procedure combines the LLM's semantic integration ability with
deterministic state management: the model decides how policy content is
organized, while the program guarantees consistency and traceability of the
update. Algorithm~\ref{alg:consolidate} summarizes the pipeline.

\begin{algorithm}[t]
\caption{Consolidation: candidate clusters $\rightarrow$ new bank version}
\label{alg:consolidate}
\begin{algorithmic}[1]
\Require candidates $\{c_i\}$, official bank $\KK_z^{(v)}$ (frozen)
\State cluster candidates: $\{\CC_g\}_{g=1}^{G} \leftarrow \Cluster(\{e(c_i)\})$
\For{each cluster $\CC_g$}
  \State retrieve candidate--bank relations $R_g$; link best relations $\ge 0.60$
  \State $P_g \leftarrow \Crud(\CC_g, \KK_z^{(v)}, R_g)$ \Comment{retry up to 3 times}
  \If{$\Valid(P_g, \KK_z^{(v)}) = 0$} \State reject cluster \EndIf
\EndFor
\Repeat
  \State build conflict graph: edge $(i,j)$ iff $W_i \cap W_j \neq \varnothing$
  \For{each connected component with $\ge 2$ plans}
    \State merge candidates, re-retrieve, re-plan jointly
  \EndFor
\Until{no conflicts}
\State run publication quality gate (payload validation, candidate support)
\State $\KK_z^{(v+1)} \leftarrow \Apply(\KK_z^{(v)},\; \bigcup_g P_g)$
\Return $\KK_z^{(v+1)}$
\end{algorithmic}
\end{algorithm}

\subsection{Validation and Frozen Evaluation}
\label{sec:method-validation}

MiM induces candidates on training failures, selects a bank version on
validation sessions, and freezes it for testing:

\begin{equation}
\label{eq:select}
v^{*} \;=\; \argmax_{v}\; \mathrm{F1}_{\mathrm{val}}\!\left(\KK^{(v)}\right),
\end{equation}

\noindent where $\mathrm{F1}_{\mathrm{val}}$ is the token-F1 on the validation
split, computed under the same frozen protocol as the test. At test time, no
reference answers are provided, no diagnosis runs, and the bank is never
updated. The base system and MiM use identical models, storage, tools, steps,
and token budgets; the only difference is the retrieval and injection of frozen
skills. To detect potential answer leakage, skills are audited for concrete
entities, dates, internal IDs, and rare strings, and evaluation is performed on
a conversation-disjoint split (6:2:2 in our setup), so that no persona, history,
or event crosses split boundaries.

Beyond maximizing task quality, MiM's long-term objective also controls the
number of policies, redundancy, and update oscillation:

\begin{equation}
\label{eq:objective}
\max_{\KK} \quad Q(\KK) \;-\; \alpha\,|\KK| \;-\; \beta\,\mathrm{Redundancy}(\KK)
\;-\; \eta\,\mathrm{Churn}(\KK) \;-\; \mu\,\mathrm{Cost}(\KK),
\end{equation}

\noindent where $Q(\KK)$ is task quality under the frozen bank. The current
method approximates this objective through discrete candidates, batch CRUD, and
version selection; it does not rely on gradient training.

\subsection{Design Rationale}
\label{sec:method-rationale}

\paragraph{Why separate construction and access skills.}
Construction and access act on different states. A construction skill changes
future memory states, so a bad update can affect many later questions; an access
skill modifies no memory and affects only how the current question obtains
evidence. Their observable traces, failure conditions, and validation methods
differ as well. A single shared bank lets rules with similar vocabulary but
different causal positions interfere---e.g., ``handle multiple temporal
versions'' may mean either preserving a validity interval during construction or
selecting the target time during access. Direction-separated banks keep repair
responsibility clear.

\paragraph{Why merge access and answer into one loop.}
Whether access should continue depends on the model's understanding of current
evidence. If an access agent produced a fixed evidence set and handed it to a
separate answer agent, the latter could not dynamically complete missing
information. MiM places retrieval and answering in the same ReAct context so
that the model can judge \textsc{Full}/\textsc{Partial}/\textsc{None} after each
observation and stop autonomously. This merge does not affect attribution: the
offline answer check still examines whether the information ultimately visible
to the running model was already sufficient.

\paragraph{Why use gold annotations only offline.}
Reference answers and evidence serve as offline supervision to judge failures
and to locate repair targets during training. The final test provides no gold
and allows no error-driven updates. This mirrors supervised training of model
parameters: gold is used to learn policies, but the frozen system must retrieve
and execute them naturally on unseen sessions.

\paragraph{Why batch candidate consolidation.}
A candidate is a local repair hypothesis induced from a single failure and does
not necessarily correspond to an independent long-term skill: multiple errors
may share one cause, and one error may already be covered by an existing skill.
Batching first preserves the independence of instance-level induction and then
applies global compression, which reduces the influence of sample order on bank
structure and provides explicit measurement objects for duplication rate, merge
rate, and version oscillation.

%=========================================================================
% 5. EXPERIMENTS — DELIBERATELY NOT WRITTEN YET
%=========================================================================
% The evaluation protocol has been defined (conversation-level 6:2:2 split of
% LoCoMo; Qwen3-8B runtime with DeepSeek-V4-Flash maintenance; frozen test with
% C/P/I LLM-as-judge and token-F1) but the experiment section is intentionally
% deferred. Planned structure:
%
% \section{Experiments}
% \subsection{Experimental Setup}
%   - data (LoCoMo, 6:2:2 conversation split), models (Qwen3-8B runtime,
%     DeepSeek-V4-Flash maintenance, MiniLM-L6-v2 embeddings), protocol
%   - baselines: Base, Global Reflection Prompt, Retrieved Failure Cases,
%     MiM-Construction, MiM-Access, MiM-Joint
%   - metrics: C/P/I LLM-as-judge, token-F1, error-type changes, skill usage
% \subsection{End-to-End Effectiveness}   (RQ1)
% \subsection{Failure Attribution Quality} (RQ2)
% \subsection{Side-Wise Ablation}          (RQ3)
% \subsection{Skill Abstraction and Reuse} (RQ4)
% \subsection{Bank Evolution}              (RQ5)
% \subsection{Cost Analysis}               (RQ7)

%=========================================================================
% 6. CONCLUSION — DELIBERATELY NOT WRITTEN YET
%=========================================================================
\section{Conclusion}
\label{sec:conclusion}

% TODO: full conclusion to be written together with the experiments section.
This paper studies how long-term memory systems can learn reusable memory
policies from their own errors. MiM decomposes downstream failures into answer
evidence utilization, current-memory access, and historical-memory construction
problems, and uses source, version, and natural search trajectories to form
auditable repair targets. Confirmed access and construction failures are
distilled into direction-separated candidate skills (access and construction),
which become an official skill bank through global relation modeling,
conflict-constrained evolution, and versioned release. The framework thereby extends the object of long-term
accumulation from ``facts about users and the world'' to ``procedural experience
about how the system itself should construct and access memory,'' and turns
policy presence, recall, and effectiveness into separately observable and
evaluable processes.

%=========================================================================
% BIBLIOGRAPHY
%=========================================================================
\begin{thebibliography}{99}

\bibitem{du2026survey}
P.~Du, ``Memory for autonomous LLM agents: Mechanisms, evaluation, and emerging
frontiers,'' arXiv:2603.07670, 2026.

\bibitem{packer2023memgpt}
C.~Packer, S.~Wooders, K.~Lin, V.~Fang, S.~G. Patil, I.~Stoica, and J.~E.
Gonzalez, ``MemGPT: Towards LLMs as operating systems,'' arXiv:2310.08560, 2023.

\bibitem{park2023generative}
J.~S. Park, J.~O'Brien, C.~J. Cai, M.~R. Morris, P.~Liang, and M.~S. Bernstein,
``Generative agents: Interactive simulacra of human behavior,'' in
\emph{Proc. UIST}, 2023.

\bibitem{zhang2025amem}
W.~Zhang, J.~Deng, X.~Zheng, and X.~Zhu, ``A-MEM: Agentic memory for LLM
agents,'' arXiv:2502.12110, 2025.

\bibitem{chattopadhyay2025mem0}
D.~Chattopadhyay, A.~Banerjee, D.~K. Harnal, W.~N. Kagitha, R.~Prasad, S.~Singh,
A.~Sharma, P.~Duttagupta, and M.~Lohia, ``Mem0: Building production-ready AI
memories for LLM applications,'' arXiv:2504.19413, 2025.

\bibitem{chen2025memoryos}
Z.~Chen, X.~Liu, X.~Li, F.~Zhou, C.~Fu, and L.~Li, ``MemoryOS: A multimodal
memory system for AI applications,'' arXiv:2506.06326, 2025.

\bibitem{vlasic2025mirix}
T.~Vla{\v{s}}i{\'c} and T.~Horvat, ``MIRIX: A multi-agent system architecture
for long-term memory,'' arXiv:2507.07957, 2025.

\bibitem{zhu2025graphiti}
Z.~Zhu, R.~R. Martinez, C.~Zhou, Z.~Li, and X.~S. Lin, ``Graphiti: Incremental
temporal knowledge graph construction for evolving world state,'' arXiv:2501.13956,
2025.

\bibitem{luo2025lightmem}
Z.~Luo, J.~He, S.~Dai, R.~Ma, B.~Wang, and Y.~Deng, ``LightMem: A lightweight
memory system for long-context LLM agents,'' in \emph{Proc. ICLR}, 2026.

\bibitem{liu2026simplemem}
J.~Liu, Y.~Su, P.~Xia, S.~Han, Z.~Zheng, C.~Xie, M.~Ding, and H.~Yao,
``SimpleMem: Efficient lifelong memory for LLM agents,'' arXiv:2601.02553, 2026.

\bibitem{van2025hindsight}
J.~van~den~Brandt, F.~Stolp, T.~de~Groot, and J.~F. Santos, ``Hindsight:
Post-hoc memory-based learning for LLM agents,'' in \emph{Proc. ACL System
Demonstrations}, 2026.

\bibitem{zhang2026memskill}
H.~Zhang, Q.~Long, J.~Bao, T.~Feng, W.~Zhang, H.~Yue, and W.~Wang, ``MemSkill:
Learning and evolving memory skills for self-evolving agents,'' arXiv:2602.02474,
2026.

\bibitem{liu2026evolvemem}
J.~Liu, X.~Ye, P.~Xia, Z.~Zheng, C.~Xie, M.~Ding, and H.~Yao, ``EvolveMem:
Self-evolving memory architecture via AutoResearch for LLM agents,''
arXiv:2605.13941, 2026.

\bibitem{memma2026}
MemMA contributors, ``MemMA: A meta-thinking approach for joint coordination of
memory construction and access in LLM agents,'' arXiv:2603.18718, 2026.

\bibitem{plugmem2026}
PlugMem contributors, ``PlugMem: A task-agnostic plug-in memory module for LLM
agents,'' arXiv:2603.03296, 2026.

\bibitem{ji2026mragent}
S.~Ji and colleagues, ``MRAgent: Active memory reconstruction with a graph-based
retrieval agent for long-term LLM agents,'' in \emph{Proc. ICML}, 2026.

\bibitem{memmachine2026}
MemMachine contributors, ``MemMachine: Ground-truth-preserving episodic memory
for LLM agents,'' arXiv:2604.04853, 2026.

\bibitem{shinn2023reflexion}
N.~Shinn, F.~Cassano, A.~Gopinath, K.~Narasimhan, and S.~Yao, ``Reflexion:
Language agents with verbal reinforcement learning,'' in \emph{Proc. NeurIPS},
2023.

\bibitem{madaan2023selfrefine}
A.~Madaan, N.~Tandon, P.~Gupta, S.~Hallinan, L.~Gao, S.~Wiegreffe, U.~Alon,
N.~Dy, S.~Gupta, A.~Borgida, K.~Sridhar, W.~Yih, and D.~Roth, ``Self-Refine:
Iterative refinement with self-feedback,'' in \emph{Proc. NeurIPS}, 2023.

\bibitem{zhong2024memorybank}
W.~Zhong, L.~Guo, Q.~Gao, H.~Ye, and Y.~Wang, ``MemoryBank: Enhancing
large language models with long-term memory,'' in \emph{Proc. AAAI}, 2024.

\bibitem{yang2024opro}
C.~Yang, X.~Wang, Y.~Lu, H.~Liu, Q.~V. Le, D.~Zhou, and X.~Chen,
``Large language models as optimizers,'' in \emph{Proc. ICLR}, 2024.

\bibitem{langmem2025}
LangChain AI, ``LangMem SDK: Long-term memory for LangGraph agents,''
\url{https://github.com/langchain-ai/langmem}, 2025.

\bibitem{wang2023voyager}
G.~Wang, Y.~Xie, Y.~Jiang, A.~Mandlekar, C.~Xiao, Y.~Zhu, L.~Fan, and A.~Anandkumar,
``Voyager: An open-ended embodied agent with large language models,''
\emph{TMLR}, 2024.

\bibitem{maharana2024locomo}
A.~Maharana, D.~Lee, S.~Tulyakov, M.~Bansal, F.~Barbieri, and Y.~Fang,
``Evaluating very long-term conversational memory of LLM agents,''
in \emph{Proc. ACL}, 2024.

\bibitem{wu2024longmemeval}
D.~Wu, C.~Zhao, Y.~Ma, and S.~Krishnan, ``LongMemEval: Benchmarking chat
assistants on long-term interactive memory,'' arXiv:2410.10813, 2024.

\end{thebibliography}

\end{document}
